\section{From context modelling to \texorpdfstring{$\cxor$}{cxor}}\label{sec:intro}

Statistical compression begins by assigning probabilities to the next symbol.  In a binary
order-$m$ context model, the $m$ most recent source bits select one of $2^m$ contexts, and the
model records what has followed that context in the past.  Prediction by partial matching (PPM),
for example, keeps symbol counts at several context orders and uses an escape mechanism when a
long context is uninformative \cite{bell1990text}.  A probability model and an entropy coder then
turn accurate probability estimates into a short representation.

Context Transformation Algorithms (CTA) denotes the broader idea of associating state with
contexts and using it to transform a source.  This paper studies one deliberately restricted
binary member, which we call $\cxor$ (CTA-XOR):
\[
  {\small
    \text{Context Transformation Algorithms (CTA)}
    \supset
    \text{$\cxor$ over $\F$}
    \supset
    \text{the historical $\W5$ instance}.}
\]

$\cxor$ can be reached from context modelling by asking how little of the architecture can
remain.  Give each context exactly one bit of state and interpret that bit as a hard prediction
$\hat x$ of the next source bit.  When the actual bit $x$ arrives, emit
\begin{equation}\label{eq:residual-intro}
  r=x\mathbin{\oplus}\hat x .
\end{equation}
Then $r=0$ means that the context predicted correctly, while $r=1$ means a prediction error.  The
result is simultaneously a context-indexed predictor and a residual transform.

The historically implemented update is as small as the state.  After observing $x$, store $x$ in
the cell for the current context:
\begin{equation}\label{eq:w5-intro}
  \W5(s,x)=x .
\end{equation}
Thus the next time the same context occurs, $\W5$ predicts that it will have the same successor as
last time.  It has one bit per context, immediate adaptation, and no counts, probabilities,
confidence, escape, or backoff.  This predictor interpretation is not an analogy added after the
fact: it is exactly the read--compare--update sequence in the earlier $\cxor$ implementation.

The residual can be more skewed than the source.  With context order $m=4$, the example shipped
with the reference $\W5$ program is
\begin{center}
\begin{tabular}{@{}r@{\ }l@{}}
input:       & \bstr{101001110111011111011111111111110}\\
$\cxor_{\W5}$: & \bstr{101001110111000011100001000000001}\\
inverse:     & \bstr{101001110111011111011111111111110}.
\end{tabular}
\end{center}
The input has $7$ zeros and $26$ ones; its image has $20$ zeros and $13$ ones.  The useful
observation is zero enrichment: repeated contextual regularity has become a concentration of
prediction successes.  The example is not evidence of compression, and no compression result is
claimed here.  $\cxor$ emits exactly one bit for every input bit and can be inverted exactly.

The paper first defines this predictor and its decoder, then makes precise which regularities
$\W5$ can expose and why their exposure is separate from both invertibility and compression.
Only after that do we vary the learning rule.  Within the one-bit $\cxor$ architecture, if a cell
may be updated using only its old value $s$ and the observed bit $x$, the update is a Boolean function
\[
  w:\mathbb F_2^2\longrightarrow\mathbb F_2,
\]
and there are exactly $2^{2^2}=16$ possibilities.  This finite design space classifies the narrow
$\cxor$ model, not the CTA family.

The construction studied here originated in earlier unpublished work by the author on the
broader CTA family.  That work existed in manuscript and implementation form, but the binary XOR
construction was never published as a separate, self-contained treatment.  The present paper
extracts that operation from the broader line of work, formalises it independently, and calls it
$\cxor$.

The analysis makes four contributions.
\begin{enumerate}
  \item It identifies the historical $\W5$ implementation as a last-observation context predictor
        followed by an exclusive-or residual.
  \item It proves that causal prediction makes every one of the sixteen transforms sequentially
        invertible; reversibility comes from the predictor/residual construction, not from a
        special feature of $\W5$.
  \item It classifies the rules by algebraic normal form.  The eight affine rules admit a compact
        per-context operator description and realise only $1$, $1+T$, or $(1+T)^{-1}$; the other
        eight have data-dependent state retention.
  \item It interprets the resulting memory horizons as adaptation behaviour.  $\W5$ replaces an
        old belief after one observation, whereas $\W6(s,x)=s\oplus x$ preserves old state.
\end{enumerate}

Three notions will remain separate throughout.  \emph{Prediction performance} is a property of a
rule, a context order, and a particular source; it can be measured by the residual zeros.
\emph{Invertibility} is structural and holds even when every prediction is poor.
\emph{Compression} requires an additional model and coder.  Confusing any two of these obscures
what $\cxor$ does.
